\section{Introduction}

A transformer is stateless between calls: everything it ``knows'' about the
current conversation must be re-supplied in the prompt. Systems that give a model
long-term memory---MemGPT, generative agents, retrieval stores---therefore treat
memory as an \emph{engineering} layer around a frozen network: text is embedded,
stored in an external database, retrieved by similarity, and pasted back into the
context window. This is effective but leaves the model itself stateless, ties
recall to the context budget, and makes ``memory'' a property of the harness
rather than of the model.

We ask a different question: can memory be a \emph{part of the model}? We add one
\emph{memory block} inside the transformer stack. It reads from a persistent,
per-model store at an early layer and writes to that store at the last block
before the head. The read injects retrieved content into the residual stream; the
write commits the late, semantically complete representation of the current
segment. With this block the model becomes \emph{stateful}: across segments it
answers questions whose evidence appears in no part of the current prompt, with
no retrieval-augmented context and no external database.

The design rests on one property. The store's key is taken at layer~0, without
position information, so a key is exactly the token embedding. Retrieval is then
an \emph{exact} token match rather than a learned soft lookup, which is what lets
a single read head recover a specific binding rather than a blurred average.
Every result in this paper depends on that layer-0 exact key; when we contextualise
the key (move it deeper, or add position) matching collapses.

This yields a clean division of labour: \textbf{the weights are stateless and the
store is the state}. A frozen model carries no memory of its own; all of it lives in
the store, which is a detachable, swappable file. We show this literally---one frozen
model with two stores answers each store's questions and only its own, and swapping the
store swaps the memory with no weight changed---and we show its boundary: a store is a
memory only for the model that wrote it, because its keys are that model's embeddings
(\S\ref{sec:transplant}).

\paragraph{Contributions.}
\begin{itemize}
\item \textbf{A stateful transformer via an in-stack memory block} (\S\ref{sec:method}):
a persistent per-model store addressed by a position-free layer-0 key; by default
every layer reads from and writes its own value into the \emph{one shared} store,
keyed by the exact token---improving the original read-early/write-late placement,
whose normalised block-0 value skip it subsumes.
\item \textbf{Cross-segment recall on real language} (\S\ref{sec:language}):
facts stated once in PG-19 prose and queried out of context are recalled at
\recall{100} through 32 facts and \recall{99} at 64, with matched controls at
chance and the store persisted across a process restart. We identify a
multi-token-key tax that caps character-level keys and show single-token
(sub-word) keys remove it, and a fact-count cliff that a warm-start curriculum
escapes.
\item \textbf{Memory beyond the context window} (\S\ref{sec:longcontext}): the
store holds far more than the attention window, answered from an empty prompt,
and is retrieved \emph{sparsely}---one to two of thousands of slots per query,
with sparse reading matching or beating a dense scan. Growing the store like a
KV cache instead of a lapping ring makes it unbounded; grown and read by a
top-$1$ lookup, recall stays flat at \recall{100} to a \emph{literal million
tokens} ($990{,}473$ slots, $\approx\!7{,}800\times$ the window) on one slot, on two
models, while the dense scan collapses to chance. A passkey (needle-in-a-haystack)
sweep of the fact's depth finds \emph{no} lost-in-the-middle but a write-order
limit---early facts persist, late facts into a saturated store are lost---traced to
a novelty write-gate and saturation, not position.
\item \textbf{A shared-KV store} (\S\ref{sec:sharedkv}): we port the three axes of
DeepSeek-V4.1-Flash's CSA2 to the store---an index-selected candidate pool
(exact token hash / LSH), cross-layer reuse of that pool, and 4-bit
quantization---each preserving recall, so the store holds $\gg$ the window, is
read in $O(1)$ slots, and costs 4 bits per element.
\item \textbf{Placement ablations} (\S\ref{sec:placement}) that independently
surface the inter-layer redundancy shared-KV schemes rely on, and locate the
fact-count cliff at a single read/write site. Read-early/write-late is in turn
improved on: reading \emph{and} writing at every layer of the one shared store,
each layer storing its own value keyed by the exact token, lifts cold recall
\recall{80.8}$\to$\recall{95.5}/\recall{98.9} (two seeds), with a residual-key
control at chance isolating the cause as per-depth values, not slot capacity.
\item \textbf{From lookup table to memory} (\S\ref{sec:memory}): a recency-biased
read that revises a re-stated fact, gated write operations that bound the store to
$O(\text{facts})$ at full recall on one shared store, and multi-hop
chain-of-thought over the store---reasoning across facts held outside the context
window---at \recall{96} (two hops) and \recall{95} (three), broken-chain control at
chance.
\item \textbf{Consolidation} (\S\ref{sec:sleep}): an offline sleep phase replays
the store into the weights with no separate decoder (the key sequence \emph{is} the
transcript), reclaiming slots while the language-model loss improves.
\item \textbf{Memory transplant} (\S\ref{sec:transplant}): a direct test of the
stateless-weights / stateful-store thesis---stores swap cleanly on the same weights
and fail across different weights---and a full chronicle of the $113$
experiments behind these results (\S\ref{sec:chronicle}), logged round by round in
\appref{app:log}.
\end{itemize}
